\documentclass[12pt, a4paper]{book}

\usepackage[T1]{fontenc}
\usepackage[utf8]{inputenc}
\usepackage[charter, cal=cmcal]{mathdesign}
\usepackage[scaled=0.92]{helvet}
\usepackage{geometry}
\geometry{a4paper, top=30mm, bottom=30mm, inner=32mm, outer=25mm, marginparwidth=18mm}

\usepackage{amsmath, amsthm}
\usepackage[version=4]{mhchem}
\usepackage{graphicx}
\usepackage{xcolor}
\usepackage{tikz}
\usetikzlibrary{arrows.meta, shapes.geometric, positioning, calc, backgrounds}
\usepackage{pgfplots}
\pgfplotsset{compat=1.18}
\usepackage{tcolorbox}
\tcbuselibrary{skins, breakable, theorems}
\usepackage{titlesec}
\usepackage{fancyhdr}
\usepackage{enumitem}
\usepackage{booktabs}
\usepackage{array}
\usepackage{parskip}
\usepackage{siunitx}
\usepackage[colorlinks=true, linkcolor=BrickRed,
            urlcolor=NavyBlue, citecolor=ForestGreen]{hyperref}

% ── Colours (Journey 5 — forest green / amber) ────────────────────────────────
\definecolor{JourneyColor}{HTML}{065F46}
\definecolor{ModuleColor}{HTML}{059669}
\definecolor{PuzzleColor}{HTML}{B45309}
\definecolor{DiscoveryColor}{HTML}{D97706}
\definecolor{WarningColor}{HTML}{991B1B}
\definecolor{ProblemColor}{HTML}{6B21A8}
\definecolor{ReflectColor}{HTML}{0D9488}
\definecolor{HistoryColor}{HTML}{78350F}

% ── Boxes ─────────────────────────────────────────────────────────────────────
\newtcolorbox{puzzlebox}[1][]{enhanced,breakable,
  colback=PuzzleColor!6,colframe=PuzzleColor,fonttitle=\bfseries\sffamily,
  title={The Puzzle},attach boxed title to top left={yshift=-2mm,xshift=6mm},
  boxed title style={colback=PuzzleColor},borderline west={3pt}{0pt}{PuzzleColor},#1}

\newtcolorbox{discoverybox}[1][]{enhanced,breakable,
  colback=DiscoveryColor!8,colframe=DiscoveryColor,fonttitle=\bfseries\sffamily,
  title={Discovery},attach boxed title to top left={yshift=-2mm,xshift=6mm},
  boxed title style={colback=DiscoveryColor},borderline west={3pt}{0pt}{DiscoveryColor},#1}

\newtcolorbox{problembox}[1][]{enhanced,breakable,
  colback=ProblemColor!5,colframe=ProblemColor,fonttitle=\bfseries\sffamily,
  title={Problems},attach boxed title to top left={yshift=-2mm,xshift=6mm},
  boxed title style={colback=ProblemColor},borderline west={3pt}{0pt}{ProblemColor},#1}

\newtcolorbox{deepbox}[1][]{enhanced,breakable,
  colback=JourneyColor!5,colframe=JourneyColor,fonttitle=\bfseries\sffamily,
  title={Deep Problems},attach boxed title to top left={yshift=-2mm,xshift=6mm},
  boxed title style={colback=JourneyColor},borderline west={4pt}{0pt}{JourneyColor},#1}

\newtcolorbox{historybox}[1][]{enhanced,breakable,
  colback=HistoryColor!7,colframe=HistoryColor!60,fonttitle=\bfseries\sffamily\small,
  title={Historical Note},attach boxed title to top left={yshift=-2mm,xshift=6mm},
  boxed title style={colback=HistoryColor!80},#1}

\newtcolorbox{reflectionbox}[1][]{enhanced,breakable,
  colback=ReflectColor!6,colframe=ReflectColor,fonttitle=\bfseries\sffamily,
  title={Reflection},attach boxed title to top left={yshift=-2mm,xshift=6mm},
  boxed title style={colback=ReflectColor},borderline west={3pt}{0pt}{ReflectColor},#1}

\newtcolorbox{namedconcept}[2][]{enhanced,
  colback=ModuleColor!6,colframe=ModuleColor,fonttitle=\bfseries\sffamily\large,
  title={#2},attach boxed title to top left={yshift=-2mm,xshift=6mm},
  boxed title style={colback=ModuleColor},#1}

% ── Section formatting ────────────────────────────────────────────────────────
\titleformat{\chapter}[display]
  {\normalfont\sffamily\huge\bfseries\color{JourneyColor}}
  {\textcolor{JourneyColor!60}{\chaptertitlename\ \thechapter}}{10pt}{\Huge}
  [\vspace{2pt}\textcolor{JourneyColor}{\titlerule[1.5pt]}]
\titleformat{\section}
  {\normalfont\sffamily\Large\bfseries\color{ModuleColor}}
  {\thesection}{1em}{}[\textcolor{ModuleColor!40}{\titlerule[0.6pt]}]
\titleformat{\subsection}
  {\normalfont\sffamily\large\bfseries\color{JourneyColor!80}}{\thesubsection}{1em}{}
\titleformat{\subsubsection}
  {\normalfont\sffamily\normalsize\bfseries}{\thesubsubsection}{1em}{}

\pagestyle{fancy}
\fancyhf{}
\fancyhead[LE]{\small\sffamily\textcolor{JourneyColor}{\leftmark}}
\fancyhead[RO]{\small\sffamily\textcolor{ModuleColor}{\rightmark}}
\fancyfoot[C]{\small\thepage}
\renewcommand{\headrulewidth}{0.4pt}

\newcommand{\newterm}[1]{\textit{\textcolor{DiscoveryColor}{#1}}}
\newcommand{\important}[1]{\textcolor{ModuleColor}{\textbf{#1}}}

\begin{document}

\setcounter{chapter}{4}
\chapter[Journey 5 --- Chemical Kinetics and Dynamics]
        {Journey Five\\[8pt]
         \textcolor{ModuleColor}{\large Chemical Kinetics and Dynamics}}
\setcounter{section}{2}

\section{Module 5.3 --- Reaction Mechanisms}
\label{sec:mod5.3}

\subsection*{Motivation}

We established in Module 5.1 that the exponents in a reaction's rate law ($R = k[\ce{A}]^x[\ce{B}]^y$) are not dictated by the coefficients of the balanced equation.
We explored in Module 5.2 how energy barriers restrict collisions to transition states.
Now, we combine these two principles to solve a deeper question: *how* do reactions actually happen?

Most chemical reactions do not occur in a single, massive collision.
A balanced equation like $\ce{2 H2 + O2 -> 2 H2O}$ represents the net summary of a complex multi-step pathway.
This pathway is the **reaction mechanism**.
This module explores how to decompose reactions into simple, individual collisions (elementary steps), identify the rate-limiting bottlenecks, and use algebraic approximations to resolve the concentration kinetics of unstable intermediates.

\subsection{Elementary Steps and rate laws}

An overall reaction is composed of a sequence of **elementary steps**, each representing a single microscopic collision event.

\begin{puzzlebox}
\textbf{The Vanishing Reactant Paradox.}

Consider the gas-phase reaction:
\[
  \ce{NO2(g) + CO(g) -> NO(g) + CO2(g)}.
\]
If this reaction occurred in a single elementary collision, the rate law would be directly dictated by its stoichiometry: $R = k[\ce{NO2}][\ce{CO}]$.
However, experimental measurements at temperatures below \SI{500}{K} reveal the rate law is:
\[
  R \;=\; k [\ce{NO2}]^2.
\]
The concentration of carbon monoxide ($\ce{CO}$) has completely no effect on the reaction speed, even though it is consumed in the reaction!
\begin{enumerate}[label=(\alph*), leftmargin=*]
  \item Why is it physically impossible for this reaction to occur in a single step matching the experimental rate law?
  \item An **elementary step** is a reaction whose rate law *is* directly dictated by its stoichiometry.
        A unimolecular step $\ce{A -> Products}$ has $R = k[\ce{A}]$. A bimolecular step $\ce{A + B -> Products}$ has $R = k[\ce{A}][\ce{B}]$.
        Explain why this direct correspondence holds for elementary steps but not for overall reactions.
  \item Suppose the reaction proceeds via a two-step mechanism:
        1. $\ce{NO2 + NO2 ->[k_1] NO3 + NO}$
        2. $\ce{NO3 + CO ->[k_2] NO2 + CO2}$
        Identify the **reaction intermediate**—a species that is produced in one step and consumed in another, and therefore does not appear in the net overall equation.
\end{enumerate}
\end{puzzlebox}

\subsection{The Rate-Determining Step}

In a multi-step pathway, the overall rate is limited by the slowest step, known as the **rate-determining step** (RDS).

\begin{explobox}
Consider a factory assembly line. If step 1 takes 10 minutes and step 2 takes 1 second, the rate of the entire line is determined entirely by the speed of step 1.
\begin{enumerate}[label=(\alph*), leftmargin=*]
  \item In the two-step mechanism for the puzzle reaction:
        1. $\ce{NO2 + NO2 ->[k_1] NO3 + NO}$ (slow)
        2. $\ce{NO3 + CO ->[k_2] NO2 + CO2}$ (fast)
        the first step is the rate-determining bottleneck.
        Write down the rate law of the rate-determining step. Show that this matches the experimental rate law of the overall reaction.
  \item Why does the concentration of $\ce{CO}$ not affect the rate? (At what point in the timeline of a single reaction path does the $\ce{CO}$ molecule get involved?)
\end{enumerate}
\end{explobox}

\subsection{Fast Pre-Equilibrium}

When the rate-determining step is preceded by a fast, reversible step, the intermediate concentration reaches a dynamic equilibrium.

\begin{explobox}
Consider the mechanism:
1. $\ce{A + B <=> [k_1][k_{-1}] I}$ (fast equilibrium)
2. $\ce{I ->[k_2] P}$ (slow, rate-determining)
where I is an intermediate.
\begin{enumerate}[label=(\alph*), leftmargin=*]
  \item The rate of production of products is dictated by the slow step: $R = k_2 [\ce{I}]$. Why is this rate law not useful as written? (Can we directly measure the concentration of an unstable intermediate in a test tube?)
  \item Since step 1 is fast and reversible, the forward and reverse rates of this step are approximately equal:
  \[
    k_1 [\ce{A}][\ce{B}] \;\approx\; k_{-1} [\ce{I}].
  \]
  Solve this equation for the intermediate concentration $[\ce{I}]$.
  \item Substitute this expression into the rate law of the rate-determining step to prove that the overall rate law is:
  \[
    R \;=\; \frac{k_2 k_1}{k_{-1}} [\ce{A}][\ce{B}].
  \]
  Identify the overall rate constant in terms of the elementary constants.
\end{enumerate}
\end{explobox}

\subsection{The Steady-State Approximation}

When we cannot assume that preceding steps reach equilibrium, we use a broader algebraic approach.

\begin{explobox}
The **Steady-State Approximation** assumes that after an initial transient startup period, the concentration of any reactive intermediate remains small and constant over time:
\[
  \frac{d[\ce{Intermediate}]}{dt} \;\approx\; 0.
\]
\problem Let us apply this to the general mechanism:
1. $\ce{A + B ->[k_1] I}$
2. $\ce{I ->[k_{-1}] A + B}$
3. $\ce{I + C ->[k_2] P}$
where I is a highly reactive intermediate.
\begin{enumerate}[label=(\alph*), leftmargin=*]
  \item Write the differential equation for the rate of change of $[\ce{I}]$, accounting for its production in step 1 and its consumption in steps 2 and 3:
  \[
    \frac{d[\ce{I}]}{dt} \;=\; k_1 [\ce{A}][\ce{B}] - k_{-1} [\ce{I}] - k_2 [\ce{I}][\ce{C}].
  \]
  \item Apply the steady-state approximation $\frac{d[\ce{I}]}{dt} = 0$, and solve for the steady-state concentration $[\ce{I}]_{\text{ss}}$.
  \item The rate of production of product is $R = k_2 [\ce{I}]_{\text{ss}} [\ce{C}]$. Substitute your expression from (b) to show that the overall rate law is:
  \[
    R \;=\; \frac{k_1 k_2 [\ce{A}][\ce{B}][\ce{C}]}{k_{-1} + k_2 [\ce{C}]}.
  \]
  \item Show that in the limit $k_{-1} \gg k_2 [\ce{C}]$, this rate law reduces to the pre-equilibrium rate law derived in Section 2.3. What is the physical meaning of this limit?
\end{enumerate}
\end{explobox}

\newpage

% ─── Section: Discovery ──────────────────────────────────────────────────────
\section{Part 3 --- Discovery}

\begin{discoverybox}
\textbf{What have we built?}

\textbf{First.} Chemical reactions occur via mechanisms composed of elementary steps. Unlike overall rate laws, the rate law of an elementary step is directly dictated by its molecular stoichiometry.

\textbf{Second.} The rate-determining step acts as a kinetic bottleneck. The overall rate law of a mechanism matches the rate law of this slowest step.

\textbf{Third.} If the RDS is preceded by a fast equilibrium, we can algebraically replace the unmeasurable intermediate concentration using the equilibrium condition $k_1[\ce{A}][\ce{B}] = k_{-1}[\ce{Intermediate}]$.

\textbf{Fourth.} The steady-state approximation ($\frac{d[\ce{Intermediate}]}{dt} = 0$) is the general algebraic method to resolve the rate laws of arbitrary mechanisms, reducing to the pre-equilibrium limit when reverse rates dominate intermediate consumption.
\end{discoverybox}

\newpage

% ─── Section: Naming Concepts ────────────────────────────────────────────────
\section{Part 4 --- What This Is Called}

\begin{namedconcept}{Reaction mechanism terms}
* **Elementary Reaction:** A reaction that occurs in a single collision event, passing through a single transition state.
* **Reaction Intermediate:** A highly reactive species produced in an early step of a mechanism and consumed in a later step.
* **Rate-Determining Step:** The slowest elementary step in a reaction mechanism that limits the overall rate.
* **Steady-State Approximation:** The assumption that the rate of change of intermediate concentrations is zero: $\frac{d[\ce{I}]}{dt} = 0$.
\end{namedconcept}

\newpage

% ─── Section: Solved Exercises ───────────────────────────────────────────────
\section{Part 5 --- Solved Exercises}

\begin{problembox}
\begin{enumerate}[label=\textbf{E\arabic*.}]
  \item \textbf{Steady-State Approximation on Ozone Decomposition.}
  
        The decomposition of ozone:
        \[
          \ce{2 O3(g) -> 3 O2(g)}
        \]
        is believed to proceed via the following two-step mechanism:
        1. $\ce{O3 <=> [k_1][k_{-1}] O2 + O}$ (fast)
        2. $\ce{O3 + O ->[k_2] 2 O2}$ (slow)
        Derive the overall rate law for the decomposition of ozone using the steady-state approximation on the oxygen atom intermediate $\ce{O}$.
        
        \textbf{Solution:}
        First, we write the differential rate equation for the intermediate $\ce{O}$.
        * $\ce{O}$ is produced by the forward reaction of step 1: rate $= k_1 [\ce{O3}]$.
        * $\ce{O}$ is consumed by the reverse reaction of step 1: rate $= k_{-1} [\ce{O2}][\ce{O}]$.
        * $\ce{O}$ is consumed by the forward reaction of step 2: rate $= k_2 [\ce{O3}][\ce{O}]$.
        
        Thus, the net rate of change of $[\ce{O}]$ is:
        \[
          \frac{d[\ce{O}]}{dt} \;=\; k_1 [\ce{O3}] - k_{-1} [\ce{O2}][\ce{O}] - k_2 [\ce{O3}][\ce{O}].
        \]
        Apply the steady-state approximation $\frac{d[\ce{O}]}{dt} \approx 0$:
        \[
          k_1 [\ce{O3}] - k_{-1} [\ce{O2}][\ce{O}] - k_2 [\ce{O3}][\ce{O}] \;=\; 0 \implies k_1 [\ce{O3}] \;=\; [\ce{O}] \left( k_{-1} [\ce{O2}] + k_2 [\ce{O3}] \right).
        \]
        Solve for the steady-state concentration $[\ce{O}]_{\text{ss}}$:
        \[
          [\ce{O}]_{\text{ss}} \;=\; \frac{k_1 [\ce{O3}]}{k_{-1} [\ce{O2}] + k_2 [\ce{O3}]}.
        \]
        The overall rate of reaction, measured by the rate of the slow step 2, is:
        \[
          R \;=\; k_2 [\ce{O3}][\ce{O}].
        \]
        Substitute the expression for $[\ce{O}]_{\text{ss}}$:
        \[
          R \;=\; \frac{k_1 k_2 [\ce{O3}]^2}{k_{-1} [\ce{O2}] + k_2 [\ce{O3}]}.
        \]
        This is the complete steady-state rate law.
        
  \item \textbf{Inhibition by Products.}
  
        For the ozone rate law derived in E1, show what the rate law reduces to under the limit $k_{-1}[\ce{O2}] \gg k_2[\ce{O3}]$. Explain the physical meaning of this limit.
        
        \textbf{Solution:}
        In the limit where the reverse rate of step 1 is much faster than the forward rate of step 2 ($k_{-1}[\ce{O2}] \gg k_2[\ce{O3}]$), we can neglect the $k_2[\ce{O3}]$ term in the denominator of the steady-state rate law:
        \[
          R \;\approx\; \frac{k_1 k_2 [\ce{O3}]^2}{k_{-1} [\ce{O2}]}.
        \]
        Defining the effective rate constant $k_{\text{eff}} = \frac{k_1 k_2}{k_{-1}}$, we obtain:
        \[
          R \;=\; k_{\text{eff}} \frac{[\ce{O3}]^2}{[\ce{O2}]}.
        \]
        Physically, this limit represents a **fast pre-equilibrium** in step 1, followed by a very slow step 2.
        Notice that the rate is inversely proportional to the concentration of the product $\ce{O2}$. This means the product $\ce{O2}$ acts as a **kinetic inhibitor**: adding oxygen gas shifts the fast pre-equilibrium backward, lowering the concentration of the reactive intermediate $\ce{O}$ and slowing the reaction.
\end{enumerate}
\end{problembox}

\newpage

% ─── Section: Problems ───────────────────────────────────────────────────────
\section{Part 6 --- Problems}

\begin{problembox}
\begin{enumerate}[label=\textbf{P\arabic*.}]
  \item \textbf{Nitrogen Dioxide and Fluorine.}
  
        The gas-phase reaction of nitrogen dioxide and fluorine:
        \[
          \ce{2 NO2(g) + F2(g) -> 2 NO2F(g)}
        \]
        has the experimental rate law $R = k[\ce{NO2}][\ce{F2}]$.
        Propose a two-step mechanism (with one slow step and one fast step) that is consistent with this rate law. Identify the intermediate.
        
  \item \textbf{Oxidation of Nitric Oxide.}
  
        The oxidation of nitric oxide:
        \[
          \ce{2 NO(g) + O2(g) -> 2 NO2(g)}
        \]
        has the experimental rate law $R = k[\ce{NO}]^2[\ce{O2}]$.
        Consider the following mechanism involving the fast pre-equilibrium formation of a dimer intermediate $\ce{N2O2}$:
        1. $\ce{2 NO <=> [k_1][k_{-1}] N2O2}$ (fast)
        2. $\ce{N2O2 + O2 ->[k_2] 2 NO2}$ (slow)
        
        \begin{enumerate}[label=(\alph*)]
          \item Derive the overall rate law for this mechanism, assuming step 1 is in fast pre-equilibrium.
          \item Does this match the experimental rate law?
          \item Explain why the rate constant of this reaction *decreases* when temperature is raised (an anomalous negative activation energy). \textit{Hint: Express the effective rate constant in terms of the equilibrium constant $K_1 = k_1/k_{-1}$, and consider how the equilibrium of the exothermic dimer formation shifts with temperature}.
        \end{enumerate}
\end{enumerate}
\end{problembox}

\newpage

% ─── Section: Deep Problems ──────────────────────────────────────────────────
\section{Part 7 --- Deep Problems}

\begin{deepbox}
\begin{enumerate}[label=\textbf{D\arabic*.}]
  \item \textbf{The General Steady-State Derivation.}
  
        Consider the general reaction mechanism:
        \[
          \ce{A + B <=> [k_1][k_{-1}] I}
        \]
        \[
          \ce{I + C ->[k_2] P}
        \]
        where I is a highly reactive intermediate.
        \begin{enumerate}[label=(\alph*)]
          \item Derive the rate law for the production of P using the steady-state approximation on intermediate I.
          \item Show that if $k_2 [\ce{C}] \gg k_{-1}$, the rate law simplifies to $R = k_1 [\ce{A}][\ce{B}]$. Explain the physical meaning of this limit. (Which step is rate-determining in this case?)
        \end{enumerate}

  \item \textbf{Chain Reactions: The Hydrogen-Bromine Reaction.}
  
        The gas-phase reaction of hydrogen and bromine:
        \[
          \ce{H2(g) + Br2(g) -> 2 HBr(g)}
        \]
        has a highly complex, non-integer rate law discovered experimentally by Max Bodenstein in 1906:
        \[
          R \;=\; \frac{k [\ce{H2}] [\ce{Br2}]^{3/2}}{[\ce{Br2}] + k' [\ce{HBr}]}.
        \]
        This reaction proceeds via a free-radical chain mechanism:
        1. **Initiation:** $\ce{Br2 ->[k_1] 2 Br^\bullet}$
        2. **Propagation:** $\ce{Br^\bullet + H2 ->[k_2] HBr + H^\bullet}$
        3. **Propagation:** $\ce{H^\bullet + Br2 ->[k_3] HBr + Br^\bullet}$
        4. **Inhibition:** $\ce{H^\bullet + HBr ->[k_4] H2 + Br^\bullet}$
        5. **Termination:** $\ce{2 Br^\bullet ->[k_5] Br2}$
        
        \begin{enumerate}[label=(\alph*)]
          \item Write down the differential equations for the rate of change of the radical intermediates $[\ce{H^\bullet}]$ and $[\ce{Br^\bullet}]$.
          \item Apply the steady-state approximation to both intermediates ($\frac{d[\ce{H^\bullet}]}{dt} = 0$ and $\frac{d[\ce{Br^\bullet}]}{dt} = 0$) to show that:
          \[
            [\ce{Br^\bullet}]_{\text{ss}} \;=\; \sqrt{\frac{k_1}{k_5} [\ce{Br2}]}, \qquad [\ce{H^\bullet}]_{\text{ss}} \;=\; \frac{k_2 [\ce{H2}] \sqrt{\dfrac{k_1}{k_5}[\ce{Br2}]}}{k_3 [\ce{Br2}] + k_4 [\ce{HBr}]}.
          \]
          \item The rate of reaction is the rate of production of the product: $R = \frac{d[\ce{HBr}]}{dt} = k_2 [\ce{Br^\bullet}][\ce{H2}] + k_3 [\ce{H^\bullet}][\ce{Br2}] - k_4 [\ce{H^\bullet}][\ce{HBr}]$.
                Substitute your steady-state expressions to derive the Bodenstein rate law. Identify $k$ and $k'$ in terms of the elementary rate constants.
        \end{enumerate}
\end{enumerate}
\end{deepbox}

\newpage

% ─── Section: Reflection ─────────────────────────────────────────────────────
\section{Part 8 --- Reflection}

\begin{reflectionbox}
\textbf{Conceptual Reflection}

Consider the following questions:
\begin{enumerate}
  \item **On the Steady State:** The steady-state approximation assumes $\frac{d[\ce{I}]}{dt} = 0$. But $[\ce{I}]$ must start at 0 and eventually return to 0 when the reactants are exhausted. Why is this approximation valid for the vast middle portion of the reaction timescale?
  \item **On Bodenstein's Rate Law:** The rate law for the hydrogen-bromine reaction contains a product term $[\ce{HBr}]$ in the denominator. Why does this imply that the product $\ce{HBr}$ acts as an inhibitor? What is the physical mechanism of this inhibition?
  \item **On Symmetries:** In the dimer pre-equilibrium model for $\ce{NO}$ oxidation, the rate constant *decreases* with temperature. Why is this "negative activation energy" not a violation of the Arrhenius law for individual elementary steps?
\end{enumerate}
\end{reflectionbox}

\medskip
\noindent\textbf{Looking ahead.}
Module~5.4 completes this journey by introducing **Catalysis**. We will study how catalysts provide alternative pathways with lower activation energies to speed up reactions. We will explore homogeneous and heterogeneous catalysis, and derive the Michaelis-Menten equations for enzyme kinetics.

\vspace{1.4cm}
\begin{center}
{\color{JourneyColor}\rule{0.5\linewidth}{0.8pt}}\\[6pt]
{\small\color{JourneyColor}\textit{End of Journey 5, Module 5.3 --- Reaction Mechanisms}}\\[4pt]
{\small\color{JourneyColor}Continue to: \textbf{Module 5.4 --- Catalysis}}\\[6pt]
{\color{JourneyColor}\rule{0.5\linewidth}{0.8pt}}
\end{center}

\end{document}
